Rapid technological advancement in computing power and data storage has enabled
the \gls{ml} field to evolve significantly. This evolution has led to the
development of versatile \gls{ml} models applied across many domains.\
\glspl{llm} have emerged as one of the most popular types of \gls{ml} models,
demonstrating various and widespread practical applications.\\

A key feature of \glspl{llm} is their ability to process natural language. The
\gls{cli} is one field where this capability can enhance user experience and
\gls{qol} for developers and regular users. It remains an essential part of
modern development processes, enabling quick and efficient task execution, and
the ability of \glspl{llm} to operate on natural language tasks can facilitate
usage of it in many ways by providing user with a much simpler interface, making
the whole interaction less demanding. This makes them valuable as assistants in
various tasks.\\

One of major well-known drawbacks of \glspl{llm}, is that for most most
solutions, it is accessed remotely on server. This introduces limitations such
as latency, dependency on network connectivity, and potential privacy concerns.
In contrast, running \glspl{llm} locally on consumer hardware offers a promising
alternative, providing low-latency, offline access and full control over the
software and information.\\

This motivates us to explore and evaluate different approaches for building a
Linux \gls{cli} assistant powered by local \glspl{llm}. The work investigates
multiple possible applications of \glspl{llm} for \gls{cli} assistance,
evaluating their practicality and performance. Through implementation and
testing, we will assess which approaches are practical for offline,
resource-constrained environments and which prove viable for real-world usage.\\
